
\section*{Общая характеристика работы}

\newcommand{\actuality}{\underline{\textbf{\actualityTXT}}}
\newcommand{\progress}{\underline{\textbf{\progressTXT}}}
\newcommand{\aim}{\underline{{\textbf\aimTXT}}}
\newcommand{\tasks}{\underline{\textbf{\tasksTXT}}}
\newcommand{\novelty}{\underline{\textbf{\noveltyTXT}}}
\newcommand{\influence}{\underline{\textbf{\influenceTXT}}}
\newcommand{\methods}{\underline{\textbf{\methodsTXT}}}
\newcommand{\defpositions}{\underline{\textbf{\defpositionsTXT}}}
\newcommand{\reliability}{\underline{\textbf{\reliabilityTXT}}}
\newcommand{\probation}{\underline{\textbf{\probationTXT}}}
\newcommand{\contribution}{\underline{\textbf{\contributionTXT}}}
\newcommand{\publications}{\underline{\textbf{\publicationsTXT}}}


{\actuality}
Интерес к изучению движения твердого тела в электрических полях стимулируется многочисленными прикладными задачами в различных областях современной техники. Подобные задачи встают при исследовании бесконтактных опор различных движущихся систем, разработке космических аппаратов~и~др.

Одной из областей, в которой неконтактные подвесы особенно актуальны, является гироскопия, а именно создание неконтактого электростатического подвеса ротора в сферическом гироскопе. Современные гироскопы с неконтактными подвесами - это сложнейшие приборы, которые вобрали в себя новейшие достижения техники. Неконтактные подвесы позволяют существенно увеличить срок службы, уменьшить трение и шум, повысить точность, увеличить рабочие скорости вращения ротора \cite{Electropribor}.

Большой интерес к разработке неконтактных подвесов проявляется как в России, так и за рубежом. Что примечательно, только три страны в мире в настоящее время способны производить электростатические гироскопы. Кроме США и Франции в их число входит и Россия \cite{Electropribor}.

С теоретической точки зрения для решения задачи движения твердого тела в электрических полях необходимо исследование совместной системы уравнений движения твердого тела около неподвижной точки и уравнений электродинамики. Масштабная работа в этом направлении проведена Ю. Г. Мартыненко в его монографии о движении твердых тел в электрических и магнитных полях \cite{Martynenko}. 


{\aim} данной работы является применение конечно элементного метода к решению задачи о сферическом роторе в электростатическом подвесе.
Начиная с простейшей одномерной модели пассивного резонансного электростатического подвеса, провести сравнение различных постановок решения связной задачи электромеханики.

Применяя методы матфизики, асимптотические методы нелинейной механики, динамики твердого тела в ходе работы аналитически оценивается поведение решения. Стоит задача получить численное решение задачи методом конечных элементов (МКЭ), сравнить с аналитическими оценками.


Для~достижения поставленной цели необходимо было решить следующие {\tasks}:
\begin{enumerate}
  \item Исследовать простейшие одноосные некотактные подвесы, провести анализ их устойчивости
  \item Исследовать и сравнить между собой способы моделирования задач электромеханики в системе конечно-элементного анализа ANSYS
  \item Разработать и исследовать, применяя метод конечных элементов,  модель пассивного одноосного подвеса
  \item Применяя методы матфизики, асимптотические методы нелинейной механики, динамики твердого тела аналитически оценить поведение решения конечно-элементной модели
  \item Разработать и исследовать, применяя метод конечных элементов,  модель сферического ротора в трехосном электростатическом подвесе 
\end{enumerate}

{\probation}
Основные результаты работы докладывались~на XIX конференции молодых ученых «Навигация и управление движением» (XIX КМУ 2017) 14-17 марта 2017 г., Санкт-Петербург, Россия.
 % Характеристика работы по структуре во введении и в автореферате не отличается (ГОСТ Р 7.0.11, пункты 5.3.1 и 9.2.1), потому её загружаем из одного и того же внешнего файла, предварительно задав форму выделения некоторым параметрам

%Диссертационная работа была выполнена при поддержке грантов ...

%\underline{\textbf{Объем и структура работы.}} Диссертация состоит из~введения, четырех глав, заключения и~приложения. Полный объем диссертации \textbf{ХХХ}~страниц текста с~\textbf{ХХ}~рисунками и~5~таблицами. Список литературы содержит \textbf{ХХX}~наименование.

%\newpage
\section*{Содержание работы}
Во \underline{\textbf{введении}} обосновывается актуальность
исследований, проводимых в~рамках данной диссертационной работы,
приводится обзор научной литературы по изучаемой проблеме,
формулируется цель, ставятся задачи работы, излагается научная новизна
и практическая значимость представляемой работы. В~последующих главах
сначала описывается общий принцип, позволяющий ... , а~потом идёт
апробация на частных примерах: ...  и~... .


\underline{\textbf{Первая глава}} посвящена ...

 картинку можно добавить так:
\begin{figure}[ht] 
  \centering
  \includegraphics [scale=0.27] {latex}
  \caption{Подпись к картинке.} 
  \label{img:latex}
\end{figure}

Формулы в строку без номера добавляются так:
\[ 
  \lambda_{T_s} = K_x\frac{d{x}}{d{T_s}}, \qquad
  \lambda_{q_s} = K_x\frac{d{x}}{d{q_s}},
\]

\underline{\textbf{Вторая глава}} посвящена исследованию 

\underline{\textbf{Третья глава}} посвящена исследованию 

В \underline{\textbf{четвертой главе}} приведено описание 

В \underline{\textbf{заключении}} приведены основные результаты работы, которые заключаются в следующем:
%% Согласно ГОСТ Р 7.0.11-2011:
%% 5.3.3 В заключении диссертации излагают итоги выполненного исследования, рекомендации, перспективы дальнейшей разработки темы.
%% 9.2.3 В заключении автореферата диссертации излагают итоги данного исследования, рекомендации и перспективы дальнейшей разработки темы.
\begin{enumerate}
  \item На основе анализа уравнений энергии одноосного подвеса была показана невозможность нахождения тела в положении устойчивого равновесия без управления напряжением.
  \item Исследована модель одноосного пассивного электростатического подвеса с переменным напряжением
  \item Описаны основные подходы при моделировании задач электромеханики в программной системе конечно-элементного анализа
  \item Произведено конечно-элементное моделирование задачи одноосного пассивного электростатического подвеса
  \item Произведен анализ и конечно-элементное моделирование задачи одноосного пассивного электростатического подвеса
  \item Произведено конечно-элементное моделирование задачи сферического трехосного пассивного электростатического подвеса
\end{enumerate}



%\newpage
При использовании пакета \verb!biblatex! список публикаций автора по теме
диссертации формируется в разделе <<\publications>>\ файла
\verb!../common/characteristic.tex!  при помощи команды \verb!\nocite! 

\ifdefmacro{\microtypesetup}{\microtypesetup{protrusion=false}}{} % не рекомендуется применять пакет микротипографики к автоматически генерируемому списку литературы
\ifnumequal{\value{bibliosel}}{0}{% Встроенная реализация с загрузкой файла через движок bibtex8
  \renewcommand{\bibname}{\large \authorbibtitle}
  \nocite{*}
  \insertbiblioauthor           % Подключаем Bib-базы
  %\insertbiblioother   % !!! bibtex не умеет работать с несколькими библиографиями !!!
}{% Реализация пакетом biblatex через движок biber
  \insertbiblioauthor           % Вывод всех работ автора
%  \insertbiblioauthorgrouped    % Вывод всех работ автора, сгруппированных по источникам
%  \insertbiblioauthorimportant  % Вывод наиболее значимых работ автора (определяется в файле characteristic во второй section)
  \insertbiblioother            % Вывод списка литературы, на которую ссылались в тексте автореферата
}
\ifdefmacro{\microtypesetup}{\microtypesetup{protrusion=true}}{}

